This report covers the final work on forecasting the state of satellite and terrestrial links for the prediction-aware multipath client built in the initial study. The task is to predict the RTT, jitter, and packet loss of a Starlink, a OneWeb, and a 5G link from 0.2\,s to 3\,s ahead, from 6\,s of 5\,Hz measurements.

\paragraph{What we found in the initial pipeline.} We reproduced the initial results and audited the data and the evaluation (Section~\ref{sec:audit}). Six issues were found and fixed. Three of them change how the initial numbers should be read:
\begin{itemize}[leftmargin=*]
  \item The Starlink parser fills missing server pings by interpolating toward the next real ping (\InterpFracSLDL\% of the downlink test targets). This leaks part of the future into the model input and makes Starlink RTT look easier to predict than it is.
  \item The earlier 95\% interval coverage was computed from the test residuals themselves and is 0.95 by construction.
  \item The horizon steps are 200\,ms, not 1\,s. The earlier ``15\,s'' horizon is 3\,s.
\end{itemize}

\paragraph{What changed in the final report.} Every method now uses causal inputs. The final forecaster keeps the initial encoder and attention but anchors the decoder on the last observation, adds causal context features, and, most importantly, predicts a mean and five quantiles for every target and step instead of a single value. Intervals are calibrated on held-out sessions with conformalized quantile regression, and three seeds are combined into an ensemble.

\paragraph{Results (test sessions, 3 seeds, against the initial model retrained on the same corrected data).}
\begin{itemize}[leftmargin=*]
  \item \textbf{Calibrated uncertainty:} the 95\% interval score of the RTT forecasts improves on all six links, by \ChgMISWorst\% to \ChgMISBest\% (mean \ChgMISMean\%), with measured coverage of \PICPMin--\PICPMax{} on unseen sessions. Jitter and loss intervals improve on every link as well.
  \item \textbf{Absolute error:} the RTT MAE drops on all six links, by \ChgMAEWorst\% to \ChgMAEBest\% (mean \ChgMAEMean\%). The gains are largest on the links with spiky RTT (5G UL \ChgMAEQUL\%, Starlink DL \ChgMAESLDL\%, OneWeb DL \ChgMAEOWDL\%).
  \item \textbf{Squared error:} RMSE and $R^2$ are essentially unchanged (RMSE \ChgRMSEBest\% to $+$\ChgRMSEWorst\%). We did not find a change that reliably lowers the RTT RMSE by more than a few percent. A tuned linear model reaches a similar RMSE skill as the initial network on four of the six links.
\end{itemize}
The improvement is therefore not in the squared error of the point forecast, which appears close to what the 6\,s window can support, but in the robustness of the point forecast to RTT spikes and in calibrated, input-dependent uncertainty. For the multipath scheduler this means a per-path risk bound that is tighter when a path is calm and widens when it is about to become unstable.
